% ###################    CONFIGURACIÓN ################### 

% Configuración de estilos generales de la página
\documentclass[12pt,
               % twocolumn
              ]{article}
              
\usepackage[paperheight=297mm,
			paperwidth=210mm,
			tmargin=28mm, 
			headheight=0mm, 
			headsep=15mm,
			textheight=240mm,
			footskip=7mm,
			textwidth=159.2mm,
			bindingoffset=10mm,
			twoside
			]{geometry} 
% -------------------------------------------------------   


% Configuración de párrafos
% -------------------------------------------------------         
\setlength{\parindent}{0em} % Tabulación de un párrafo nuevo
\setlength{\parskip}{1em}   % Distancia en blanco entre un párrafo y otro.
\setlength{\columnsep}{1em} % Margen izquierdo de columnas
\renewcommand{\baselinestretch}{1.5} % Interlineado. Espacio entre dos renglones.
% ------------------------------------------------------- 


% Paquetes de idioma
\usepackage[utf8]{inputenc}
\usepackage[spanish,english,es-tabla]{babel} % English para el abstract
% ------------------------------------------------------- 

% Paquetes para imprimir código fuente
\usepackage[cache=false]{minted}

\usepackage{subcaption}

\usepackage[table,xcdraw]{xcolor} % Color de tabla

\usepackage{movie15}

% -------------------------------------------------------  

% Paquetes generales
\usepackage{enumitem} % Para configurar enumeraciones
%\usepackage[usenames,dvipsnames]{xcolor}    % Para colores en el texto

% Lorem ipsum
\usepackage[pangram]{blindtext}

\usepackage{scrextend} % Para entorno{addmargin} (dar margen a un párrafo...)

% Lorem ipsum
\usepackage{lipsum}

\usepackage{fancyhdr}

\usepackage{float} % para controlar la situación de los entornos flotantes

\usepackage{graphicx} % figuras

% Definimos un nuevo color
\definecolor{webred}{rgb}{0.5, 0, 0}   % less intense red

\usepackage[bookmarks=true,
			bookmarksnumbered=false, % true means bookmarks in 
									 % left window are numbered                         
			bookmarksopen=false,     % true means only level 1
									 % are displayed.
			colorlinks=true,
			linkcolor=webred,
			%linkcolor=OliveGreen,
			citecolor=webred,
			urlcolor=cyan]{hyperref} % Para hiperenlaces

% Definimos una nueva orden para pintar una línea cuando escribamos \HRule
\newcommand{\HRule}{\rule{\linewidth}{1mm}}




% ###################   FIN CONFIGURACIÓN ################### 

% ###################	INICIO DEL DOCUMENTO ###################

\begin{document}

\thispagestyle{empty}

	\begin{figure}[h]
		\centering
		\includegraphics[width=0.7\textwidth]{images/logoUGR.jpg}
		\label{imagen2}
	\end{figure}

	\begin{center}
		\Large{TRABAJO FIN DE MÁSTER} \\
		\large{INGENIERÍA INFORMÁTICA}
	\end{center}
	
	\vspace{0.5cm}
	
	\hrule
	
	\begin{center}
		\huge{\textbf{Sistema de videovigilancia con una Raspberry PI}}
	\end{center}
	
	\hrule
	
	\vspace{0.5cm}
	
	\begin{center}
		\textbf{Autor} \\
		Jonathan Martín Valera (alumno)\\
		\textbf{Directores}\\
		Juan Julián Merelo Guervós (tutor)
	\end{center}
	
	\vspace{0.6cm}
	
	\begin{figure}[h]
		\centering
		\includegraphics[width=0.3\textwidth]{images/etsiit_logo.png}\\[0.1cm]
	\end{figure}
	
	\begin{center}
	\textsc{Escuela Técnica Superior de Ingenierías Informática y de Telecomunicación}\\
	\textsc{---}\\
	Granada, septiembre de 2019
	\end{center}
	
	\newpage

\thispagestyle{empty}

\section*{Agradecimientos}

Después de todo este periodo de formación académica, puedo mirar hacia atrás y hacer una retrospectiva de todo lo sucedido en estos años.

Ha sido una etapa bastante dura e intensiva; una etapa donde se ha tenido que superar grandes obstáculos y emplear muchísima dedicación y esfuerzo. Puedo decir con satisfacción que ha valido la pena. Considero que en esta etapa he adquirido una gran madurez mental, en la que no solo he aprendido algunos conceptos específicos, sino que que puedo decir que he aprendido a pensar y a ser autocrítico con el trabajo.

Gracias al grado y máster de ingeniería informática, veo la realidad de otra forma, con capacidad de poder realizar cualquier reto que me pueda proponer. Si tuviera que resumir en una frase todo lo aprendido en estos años, sería: \textit{`No esperes a que nadie te diga como hacer las cosas, simplemente investiga, experimenta y saca tus propias conclusiones'}.

Este trabajo fin de máster es un ejemplo de ello, una idea hipotética, en la que he investigado, aprendido lo necesario para llevarla a cabo e implementado con éxito.

Quiero agradecer a todos los profesores que he tenido a lo largo de estos años, con especial atención a mi tutor Juan Julián Merelo Guervós, que me ha enseñado en este último año a investigar y conocer una gran cantidad de herramientas actuales que me están siendo de gran utilidad en mi día a día laboral.

Finalmente, agradecer a mis padres la ayuda que, a su forma cada uno, me han proporcionado
en todo este periodo, a mis tíos por apoyarme, motivarme y confiar siempre en mí, y
especialmente a mi novia Nerea, que me ha estado animando todo este tiempo y dando apoyo
moral además de ser fuente de inspiración y superación en mí.

\newpage



\thispagestyle{empty}

\vspace*{\fill} % Para alinear el contenido de la paǵina verticalmente.

\begin{center}
	\textbf{Sistema de videovigilancia con una Raspberry PI}\\
	
	\vspace{0.3cm}
	
	Jonathan Martín Valera (alumno)	\\
	
	\footnotesize{Estudiante de Informática en Escuela Técnica Superior Informática y Telecomunicación, Universidad de Granada, 18071, Granada, España} \\
	
\end{center}

\selectlanguage{spanish}

\begin{abstract}

\noindent \textbf{Palabras clave}: videovigilancia, sistema de seguridad, raspberry PI. 

\vspace{0.2cm}

\noindent Este proyecto se basa en el desarrollo e implementación de un sistema de videovigilancia de bajo coste, haciendo uso de una raspberry PI y sus componentes. La idea es tener una raspberry PI equipada con una cámara y un sensor de movimiento para poder utilizarla como sistema de seguridad, controlada a través de un dispositivo móvil usando la aplicación de Telegram. 

\noindent La aplicación consta de una serie de modos (manual, automático y streaming) con los que se puede vigilar la zona que desees, y enviar alertas automáticas a tu móvil en el caso de detectar alguna actividad sospechosa, capturar imágenes o grabar vídeos de forma manual, o visualizar la retransmisión en directo de la visualización de la cámara, y todo esto se puede controlar desde cualquier localización.

\end{abstract}

\selectlanguage{english} 

\begin{abstract}

\noindent \textbf{Keywords}: video surveillance, security system, raspberry PI.

\vspace{0.2cm}

\noindent This project is based on the development and implementation of a low-cost video surveillance system, using a raspberry PI and its components. The idea is to have a PI raspberry equipped with a camera and a motion sensor to be used as a security system, controlled through a mobile device using the Telegram application. 	

\noindent The application consists of a series of modes (manual, automatic and streaming) with which you can monitor the area you want, and send automatic alerts to your mobile in the event of detecting any suspicious activity, capture images or record videos manually, or view the live broadcast of the camera display, and all this can be controlled from any location.

\end{abstract}

\vspace*{\fill} % Para alinear el contenido de la paǵina verticalmente.

\selectlanguage{spanish}

\newpage


\fancypagestyle{miEstilo1}{
   \lhead{1. Introducción}
   %\chead{1. Introducción}
   \rhead{Página \thepage}
   \lfoot{}
   \cfoot{}
   \rfoot{}
}

\pagestyle{miEstilo1}

\section{Introducción}

\subsection{Motivación}

El número de robos y hurtos ha ido creciendo en estos últimos años, y por ello, es necesario la aplicación de medidas preventivas y correctoras que intenten reducir este número.

Hoy en día vivimos en una sociedad informatizada y tecnológica, donde cada uno tiene la capacidad de poder encontrar casi cualquier información allá donde esté, gracias a los dispositivos móviles y a la redes de comunicaciones que usa internet.

¿Quién no lleva un smartphone con conexión de datos en su bolsillo? En España el 85\% de la población hace uso de su smartphone a diario. Estos dispositivos móviles tienen una gran cantidad de usos y aplicaciones, y se han convertido en una herramienta imprescindible en nuestro día a día.

Uniendo tecnología y aplicación de medidas de seguridad basada en videovigilancia, surge la idea de este proyecto; la idea de construir un sistema de videovigilancia de bajo coste que sea fácil de usar, mantener, y lo más importante, que pueda ser controlado desde cualquier parte.

Partiendo de esta idea, se ha investigado y construido una aplicación que cumple con estas expectativas y más, ya que el uso que se le quiera dar puede ir más allá del concepto de videovigilancia para la seguridad. Por ejemplo, se puede utilizar para controlar entradas y salidas en una zona, control parental \ldots, aunque eso sí, dentro de un uso permitido y responsable.

Otro de los principales aspectos de este proyecto, es que su uso no sea de forma privada, es decir, que pueda ser usado por todas las personas que lo deseen. Por este motivo, se ha intentado reducir todo lo posible el coste del hardware, y se ha liberado todo el código fuente, junto con las instrucciones necesarias para la instalación y despliegue de la aplicación.

Toda la información relacionada se puede consultar en el repositorio de Github \cite{ref1}.

\subsection{Objetivos}

El objetivo principal de este proyecto es el de poder construir un sistema de seguridad basado en la videovigilancia, de bajo coste, y accesible a todo el mundo.

\textbf{Objetivos generales}

\vspace{-0.3cm}

\begin{itemize}
	\item Rebajar el coste de los sistemas de videovigilancia habituales.
	\item Alertar a un usuario ante un evento de movimiento generado dentro de una zona controlada.
	\item Almacenar y obtener pruebas (fotos y/o vídeo) tras la detección de cualquier intruso.
	\item Monitorizar y controlar el estado de un entorno.
	\item Aumentar el nivel de seguridad de cualquier tipo de entorno.	
\end{itemize}

\textbf{Objetivos específicos}

\vspace{-0.3cm}

\begin{itemize}
	\item Proporcionar un sistema accesible para todo el mundo y fácil de usar.
	\item Permitir el acceso y gestión del sistema a través de un dispositivo móvil.
	\item Investigar y conocer el uso de los bots de Telegram como medio de interacción entre el usuario y la aplicación back-end desarrollada en este proyecto.
	\item Estudiar el funcionamiento de una Raspberry PI y sus principales componentes para su aplicación en el ámbito del proyecto.
	\item Conocer, diseñar y hacer uso de una arquitectura basada en microservicios que permita el uso de servicios independientes, ágiles y escalables.
	\item Aprender a utilizar mecanismos de aprendizaje automático para poder filtrar y evitar falsos positivos en las alertas generadas.

\end{itemize}

\newpage

\subsection{Estructura del documento}

Este documento está estructurado de la siguiente forma:

\textbf{1. Introducción}: Motivos y objetivos por los cuales se ha desarrollado este proyecto.

\textbf{2. Planificación del proyecto}: Muestra información detallada acerca de las diferentes etapas del proyecto, junto con la planificación prevista.

\textbf{3. Análisis de mercado}: Estudio comparativo acerca de softwares similares al propuesto en este proyecto.

\textbf{4. Tecnologías y herramientas utilizadas}: Breve descripción sobre las tecnologías y herramientas utilizadas para desarrollar este proyecto.

\textbf{5. Descripción de la aplicación: SIVIRA}: Descripción sobre los componentes de la aplicación, la interacción entre ellos, medidas de seguridad aplicadas, interfaz de usuario y descripción de la estructura de archivos y directorios de la aplicación.

\textbf{6. Presupuesto}: Estimación de costes del desarrollo del proyecto junto con una previsión de futuros beneficios en el mercado laboral.

\textbf{7. Conclusiones y trabajos futuros}: Conclusiones tras finalizar el desarrollo del proyecto y posibles mejoras a desarrollar en versiones posteriores.

\textbf{ANEXO}: Información útil sobre la aplicación SIVIRA: guía de instalaciones de componentes, instalación y despliegue de la aplicación y finalmente una guía de usuario.

\newpage



%

\fancypagestyle{miEstilo2}{
   \lhead{2. Planificación}
   %\chead{1. Introducción}
   \rhead{Página \thepage}
   \lfoot{}
   \cfoot{}
   \rfoot{}
}

\pagestyle{miEstilo2}

\section{Planificación}

Como en todo proyecto, realizar una buena planificación es la clave para no fracasar en el desarrollo del proyecto. Por ello, se ha realizado una planificación detallada del proyecto, utilizando el concepto de \textit{iteraciones} y \textit{sprints} para ir distribuyendo el trabajo a lo largo de las semanas, durante las cuales se han ido desarrollando una serie de tareas que cumplen un objetivo específico.

Estos conceptos surgen de los llamados métodos de desarrollo ágiles, en los cuales se intenta descomponer una aplicación o proyecto por funcionalidades, y el objetivo es ir desarrollando y testeando cada una de forma independientes para poder integrarla con el resto.

A continuación se muestra la planificación que se ha previsto para el desarrollo del proyecto.


\subsection{Fase preparatoria}

En esta fase se comenzará el desarrollo del proyecto. Tiene un esfuerzo estimado de X horas. Se realizará un estudio de viabilidad de la idea del proyecto, así como un estudio de mercado, una investigación acerca de las posibles herramientas a usar \ldots.

El objetivo de esta primera fase es tener claro que el proyecto es viable tanto en recursos, tiempo y presupuesto, para que posteriormente comience su desarrollo e implementación.

\large{\textbf{Iteración 1}: Análisis y estudio de la aplicación.}
\hrule

\vspace{0.3cm}

\normalsize

En esta iteración se describirán los principales objetivos de la aplicación, se estudiará sus posibles casos de uso y se elaborará un plan de iteraciones para planificar el desarrollo del proyecto a lo largo del tiempo dispuesto.

Esta iteración se divide en los siguientes sprints:

\textbf{Sprint 1}: Descripción de la aplicación



%\begin{center}
%\begin{tabular}{ | m{5em} | m{1cm}| m{1cm} | } 
%\hline
%cell1 dummy text dummy text dummy text& cell2 & cell3 \\ 
%\hline
%cell1 dummy text dummy text dummy text & cell5 & cell6 \\ 
%\hline
%cell7 & cell8 & cell9 \\ 
%\hline
%\end{tabular}
%\end{center}




\begin{table}[]
\begin{tabular}{|p{4cm}|p{7cm}|p{1.5cm}|p{2.5cm}|}
\hline
\rowcolor[HTML]{000000} 
{\color[HTML]{FFFFFF} Product backlog} & {\color[HTML]{FFFFFF} Descripción}                                  & {\color[HTML]{FFFFFF} Semana} & {\color[HTML]{FFFFFF} T.Previsión} \\ \hline
Objetivos                              & Descripción de los objetivos                                        & 1                             &                                        \\ \hline
Funcionalidades                        & Descripción de las funcionalidades                                  & 1                             &                                        \\ \hline
Motivaciones                           & Motivaciones personales acerca de la temática del proyecto          & 1                             &                                        \\ \hline
Usuarios y escenarios                  & Descripción de los posibles casos de uso de la aplicación propuesta & 1                             &                                        \\ \hline
\end{tabular}
\end{table}






\newpage

\pagestyle{empty}

\begin{thebibliography}{99}

	\bibitem{ref1} Jonathan Martín Valera, TFM Security system PI github repository, disponible en \url{https://github.com/jmv74211/TFM_security_system_PI}

	\bibitem{ref2} Gus, 26 Jul 2019, Build a Raspberry Pi Security Camera Network,	disponible en \url{https://pimylifeup.com/raspberry-pi-security-camera/}
	
	\bibitem{ref3} ccrisan, 3 Sep 2019, A Video Surveillance OS For Single-board Computers ,	disponible en \url{https://github.com/ccrisan/motioneyeos}
	
	\bibitem{ref4} scavix, instructables.com, Raspberry Pi As Low-cost HD Surveillance Camera disponible en \url{https://www.instructables.com/id/Raspberry-Pi-as-low-cost-HD-surveillance-camera/}
	
	\bibitem{ref5} Motion-project, 19 Ago 2019, Camera video monitor, disponible en \url{https://github.com/Motion-Project/motion}
	
	\bibitem{ref6} Humberto Higinio, 14 Oct 2018, Sistema de Seguridad Raspberry Pi-Sensor de movimiento y cámara HD con envío de imágenes a correo, disponible en \url{https://www.youtube.com/watch?v=rK6uLwMRtIs}
	
	\bibitem{ref7} Google, DrawIO, diseño y diagramas de aplicaciones open source, disponible en \url{https://www.draw.io/}
	
	\bibitem{ref8} Visual paradigm, diseño de diagramas UML, disponible en \url{https://www.visual-paradigm.com/}
	
	\bibitem{ref9} Israel Alcázar, 03 Jun de 2014,  Introducción a Git y Github, disponible en \url{https://desarrolloweb.com/articulos/introduccion-git-github.html}
	
	\bibitem{ref10} Git, Sistema de control de versiones distribuido, multiplataforma y de código abierto, disponible en \url{https://git-scm.com/}
	
	\bibitem{ref11} Github, Servicio de alojamiento para el desarrollo de software utilizando Git , disponible en \url{https://git-scm.com/}
	
	\bibitem{ref12} PiCamera documentación, Interfaz de cámara Raspberry PI para python.
	
	\bibitem{ref13} pyTelegramBotAPI, Bibliotecas de funciones para la conexión con la API de telegram, disponible en \url{https://github.com/eternnoir/pyTelegramBotAPI}.
	
	\bibitem{ref14} Flask, framework para el despliegue ágil de aplicaciones sencillas, disponible en \url{https://flask.palletsprojects.com/en/1.1.x/}.
	
	\bibitem{ref15} Celery, Cola de tareas distribuidas, disponible en \url{http://www.celeryproject.org/}.
	
	\bibitem{ref16} RPi.GPIO, Biblioteca para controlar los pines GPIO de la Raspberry PI, disponible en \url{https://pypi.org/project/RPi.GPIO/}.
	
	\bibitem{ref17} Tensorflow, An end-to-end open source machine learning platform, disponible en \url{https://www.tensorflow.org/}.
	
	\bibitem{ref18} Requests, An end-to-end open source machine learning platform, disponible en \url{https://www.tensorflow.org/}.
	
	\bibitem{ref19} Latex-project, An introduction to LaTeX, disponible en \url{https://www.latex-project.org/about/}.

	\bibitem{ref20} TeX Live, 06 Jul de 2019, Introduction to TeX Live, disponible en \url{https://www.tug.org/texlive/}.
	
	\bibitem{ref21} TeXstudio.org, Welvome to TeXstudio, disponible en \url{https://www.texstudio.org/}.

	\bibitem{ref22} Jetbrains, Pycharm, The Python IDE
	for Professional Developers, disponible en \url{https://www.jetbrains.com/pycharm/}.
	
	\bibitem{ref23} RabbitMQ, Most widely deployed open source message broker, disponible en \url{https://www.rabbitmq.com/}.
	
	\bibitem{ref24} Telegram bot, Bots: An introduction for developers, disponible en \url{https://core.telegram.org/bots}.
	
	\bibitem{ref25} Olprod, 13 Nov 2018, Estilo de arquitectura de microservicios, disponible en \url{https://docs.microsoft.com/es-es/azure/architecture/guide/architecture-styles/microservices}.

	\bibitem{ref26} Chakray, 18 Jul 2018, DevOps: Arquitectura monolítica vs Microservicios, disponible en \url{https://www.chakray.com/es/devops-arquitectura-monolitica-vs-microservicios/}
	
	\bibitem{ref27} Tensorflow models ,Tensorflow detection model zoo, disponible en \url{https://github.com/tensorflow/models/blob/master/research/object_detection/g3doc/detection_model_zoo.md}
	
	\bibitem{ref28} MDN web docs, The WebSocket API (WebSockets), disponible en \url{https://developer.mozilla.org/en-US/docs/Web/API/WebSockets_API}
	
	\bibitem{ref29} waveform Pistreaming repository, A little demo of streaming the Pi's camera to web browsers, disponible en \url{https://github.com/waveform80/pistreaming}.
	
	\bibitem{ref30} Margaret Rouse, RESTful API, disponible en \url{https://searchapparchitecture.techtarget.com/definition/RESTful-API}
	
	\bibitem{ref31} Core telegram , Telegram APIs and Bot API, disponible en \url{https://core.telegram.org/}
	
	\bibitem{ref32} eternnoir pyTelegramBotAPI, A simple, but extensible Python implementation for the Telegram Bot API, disponible en \url{https://github.com/eternnoir/pyTelegramBotAPI}
	
	
	\bibitem{refx1} Core telegram, New Inline Keyboards, disponible en \url{https://core.telegram.org/bots/2-0-intro}
	
\end{thebibliography}



\end{document}

% ###################	FIN DEL DOCUMENTO ###################